\section{Named objects and signed-spin angles}
\label{ten:named}
This section proves a fixed-moment estimate that allows a representation to be described by two signed-spin partitions and a set of distinct named positions. The induction keeps the complete dependence on the number of names. The negative density term in their budget pays for the number of ways to place them in a rectangle.

Proposition~\ref{ten:named:main} below also follows directly from Theorem~\ref{s:spin-moment}; the short deduction is given immediately after its statement. The development in this section supplies a separate proof: the signed-spin setup, crowded-cell sparse estimate, and coefficient-integral angle estimate prepare the trace induction, which closes the same named-object budget. We retain both routes because the second provides these distinct methods.
\smallskip
The notation below is local; the sweep and trace conventions remain those of Section~\ref{sec:prelim}.

\subsection{A budget for named objects}
We use representation theory and a trace induction for a whole sweep of switches. Two features of the induction are:

\begin{itemize}
\item For a trace test one can use a mixture of named objects and tensor spins, not just the regular representation or just spins. The named objects will have to return to the very same cells of a rectangular array. There is an angle saving on the remaining spin space.
\item It suffices to propagate bounds which can be quite large, but must be small compared with the dimension of an irrep. It helps in doing this to give the named objects a budget per object which \textbf{decreases} at low density (but is clipped at zero).
\end{itemize}

We use partitions with irreps \([\lambda]\), dimensions \(D_\lambda\), and transposed diagrams \(\lambda^{\rm t}\). For partitions of separate sizes we denote induction from the Young subgroup by a product notation such as \(\mu*\xi\), also used with more factors; membership denotes occurrence. Transposing all diagrams preserves occurrence. Throughout, diagrams of size zero and their representations are allowed.

If \(z\) is a list of nonnegative numbers with total \(u\), put
\[
 \mathcal S(z)=\sum_{j:z_j>0} z_j\log(u/z_j);
\]
we allow empty lists and use zero when the total is zero. With two partitions as arguments this takes the concatenated list of row lengths, not the entropy of each separately.

We bound matrices for a single sweep in terms of a spin and named-object budget. Define
\begin{equation}
 W_n(\lambda)=D_\lambda\operatorname{Tr}|B_n([\lambda])|^{2h},
 \label{ten:named:eq:7}
\end{equation}
where we will choose an absolute integer \(h\), possibly very large. Here \(|\cdot|\) takes the operator absolute value.

Fix the numerical budget parameters
\[
 \delta=.00001,\qquad
 \beta_n=.18-\frac{.04}{1+\log n}.
\]
Write \([x]_+=\max\{0,x\}\) and define
\begin{equation}
 R_n(z)= z[\,\delta\log n-\log(n/z)-C_0\,]_+
 \qquad(0\le z\le n),
 \label{ten:named:eq:8}
\end{equation}
with \(R_n(0)=0\) and a sufficiently large absolute \(C_0\) to be fixed.

\begin{proposition}[Named-object moment bound]
\label{ten:named:main}
For the constants fixed above, there are absolute choices of \(C_0\) and
an integer \(h\ge1\) such that, for every dyadic \(n\ge2\), every
\(\lambda\vdash n\), and every occurrence shown below with \(m=|\tau|\),
\begin{equation}
 \lambda\ \text{in}\ \mu*\nu^{\rm t}*\tau
 \quad\Longrightarrow\quad
 W_n(\lambda)
 \le
 \exp\{ \beta_n \mathcal S(\mu,\nu)+ R_n(m)\}.
 \label{ten:named:eq:9}
\end{equation}
\end{proposition}

\begin{proof}[Deduction from the every-occurrence bound]\label{s:named-specialization}
The occurrence in~\eqref{ten:named:eq:9} is exactly~\eqref{ten:spin:eq:1} by Frobenius reciprocity, and $\mathcal S(\mu,\nu)=F(\mu,\nu)$. Choose $\bar\eta<.14$ in Theorem~\ref{s:spin-moment} and then $\bar\kappa\le\delta/2$, also satisfying that theorem's parameter restrictions. For $\log n\ge2C_0/\delta$,
\[
 \eta_d\le\beta_n,\qquad
 [\kappa_d\log n-\log(n/m)]_+
 \le[\delta\log n-\log(n/m)-C_0]_+.
\]
The inequality also holds at $m=0$ with the stated conventions. Thus the theorem gives the desired result at every sufficiently large dyadic size with $h\ge r$. The finitely many smaller nontrivial Fourier blocks satisfy a strict norm gap by Lemma~\ref{lem:finitegap}; increasing the fixed $h$ makes their weighted moments at most one. The trivial block already has moment one, and the target budget is nonnegative. This proves the stated constants and quantifiers.
\end{proof}

\subsection{Signed-spin representations}
For the separate proof, we recall the representation facts used below:

\begin{itemize}
\item We use the regular representation/Fourier formula (in particular Plancherel), Schur-Weyl duality with polynomial \(GL\) irreps, the Young branching, Pieri and Littlewood-Richardson rules, the hook and Weyl dimension formulas, and unitary Schur orthogonality. Recall that if \(\lambda\) occurs in a two-factor product, it contains each factor's diagram. When a diagram contains another as upper left subdiagram, the bigger occurs in the product of the smaller with \emph{some} diagram of the remaining size, by branching.
\item We use the polynomial irreps of \(GL(q)\) also as \(U(q)\) types. Their dimensions, on degrees bounded by \(n\), are \(\exp(O(q^2\log(n+2)))\) when \(q\le n\). The matrix of a positive diagonal with descending spectrum \(y_1,\ldots,y_q\) on a type with partition \(\alpha\) has norm \(\prod_j y_j^{\alpha_j}\). We can use the Hilbert spaces as in unitary Schur-Weyl duality; this formula follows, for instance, by highest weight and dominance of the weights by the highest weight. A spectrum including zeros follows by continuity.
\end{itemize}

Here are two useful descriptions of spin types. Take a space
\[
 V=\mathbb C^q_+\ \oplus\ \mathbb C^q_-
\]
of plus and minus spins, and take \(G=U(q)\times U(q)\), acting by its two defining representations. Let position permutations act on tensor slots by graded permutations: the minus slots incur signs, that is, an interchange of two minus spins has an additional \(-1\). This commutes with \(G\). On \(V^{\otimes u}\), the \(G\)-types are indexed by pairs \(w=(\alpha,\gamma)\) of partitions with at most \(q\) parts each, with total size \(u\). We use the full isotypic spaces. As a module over the symmetric group, the space of such a type consists of copies of \(\alpha*\gamma^{\rm t}\), the number of copies being the dimension of the \(G\)-irrep. Indeed, one takes plus and minus counts, uses ordinary Schur-Weyl on both, with the sign twist on the latter positions, and induces.

For \(u\le n\), any \(\eta\) in \(\alpha*\gamma^{\rm t}\) satisfies
\begin{equation}
 \log D_\eta\ \ge\ \mathcal S(\alpha,\gamma)-C q^2\log(n+2)
 \qquad (q\le n).
 \label{ten:named:eq:1}
\end{equation}
Here and in coarse errors, \(C\) denotes an absolute constant that may change. To see this, one can reach \(\eta\) from \(\alpha\) by adding at most \(q\) vertical strips (since \(\gamma^{\rm t}\) occurs in a product of that many column diagrams). And similarly one can reach \(\eta^{\rm t}\) from \(\gamma\). In particular \(\eta\) lives in the \(q,q\) hook, and its long row and column lengths agree with those of \(\alpha,\gamma\) within \(O(q)\) each. In the dimension formula, use at most \(O(q^2)\) factors bounded by \(2n\) for the square region, and in a row protruding to the right of it bound hook lengths by the remaining row counts (including the current box), each increased by \(q\); similarly below the square by columns. Their factorial products are bounded by the products of row factorials of \(\alpha\) and of \(\gamma\), times \(\exp(C q^2\log(n+2))\). This and the factorial estimate for entropy give \eqref{ten:named:eq:1}. In the other direction, total dimension of \(\alpha*\gamma^{\rm t}\) is at most \(\exp(\mathcal S(\alpha,\gamma))\), by the dimension formula for induction and the elementary factorial-quotient bounds for the two Specht dimensions and for multinomial coefficients. Thus multiplicities in a spin type, including \(G\)-dimensions, cost at most \(\exp(C q^2\log(n+2))\).

\subsection{Sparse path cancellation}
We give first a separate lemma for those representations for which an estimate using few cards will suffice. Take any fixed \(c\in(0,1)\). For all sufficiently large \(n\), if \(1\le k=n-\lambda_1\le n^{1-c}\), we have
\begin{equation}
 \|B_n([\lambda])\|\ \le\ \exp\{-c' k\log(n/k)\}
 \label{ten:named:eq:2}
\end{equation}
with some fixed \(c'>0\). Constants in this lemma can use \(c\).

Use the action on \(k\)-tuples of distinct positions. The irrep occurs here but not with fewer positions, by branching. For tuples \(x,y\), for \(J\subseteq\{1,\ldots,k\}\), write \(K_J(x,y)\) for \(n^{|J|}\) times the probability that the sweep takes \(x_J\) to \(y_J\), and put
\[
 L(x,y)=\sum_J(-1)^{k-|J|} K_J(x,y).
\]
For fixed \(x_i,y_i\) a path in the switch network is specified: each refreshed bit has its required final value. If one of the paths in \(x,y\) does not even share a switch with any other, \(L=0\). Indeed it costs \(1/n\) in all the relevant sweep probabilities, with independent switches. On the other hand, \(L\), as a kernel with multiplier \(n^{-k}\) on the distinct tuples, gives the action of the sweep on \([\lambda]\) after compression to it: the proper-subset terms vanish, as their ranges or co-ranges involve fewer entries. Thus we will bound the op norm of the matrix \(n^{-k}L\). If \(k=1\) the lemma is immediate, so take \(k\ge 2\).

We use the dyadic cells on starts \(x\) in which bits not yet swept are fixed; on the end positions \(y\) use the corresponding cells for the reverse order. Each is a family nested as a binary tree. Put
\[
 p_*=(k/n)^{1/10}.
\]
Call a position in the tuple crowded if it lies in a cell containing at least two of the entries and with entry density at least \(p_*\). A tuple is bad if at least \(k/4\) entries are crowded. For a noncrowded start entry \(i\) there is a useful first-contact bound:
\begin{equation}
 \sum_{j=1}^d 2^{-(j-1)}
 \#\left\{\begin{array}{l}
       \text{other starts whose bits remaining}\\
       \text{after step }j\text{ agree with entry }i
      \end{array}\right\}\ \le\ 2 d p_*.
 \label{ten:named:eq:3}
\end{equation}
Indeed a contributing cell has size \(2^j\).

Think of \(n^{-k}K_J\) in a row as generating paths: actual shuffled paths for the entries in \(J\), and independent uniform ends (with the corresponding paths) for the others, even allowing nondistinct ends until restricting the matrix. For any order of inspecting paths, the conditional probability the next path of a noncrowded start will hit any earlier paths, meaning share any switch, is bounded by \eqref{ten:named:eq:3}. Until its first hit it makes free uniform choices. For shuffled cards, the earlier shuffled paths condition only their own visited switches. And at layer \(j\) a possible hit with any fixed earlier path requires the first \(j-1\) generated bits to agree with that path at the given switch, as well as agreement on the unswept bits other than bit \(j\), just as counted in \eqref{ten:named:eq:3}.

If \(x\) is not bad, the row probability of all vertices of the path contact graph being nonisolated costs \(\exp(-\Omega(k\log(n/k)))\), uniformly for each \(J\). For clarity, inspect in an independently random order. Given this graph without isolated vertices, at least \(3k/4\) noncrowded starts each have probability at least \(1/2\) of an earlier neighbor, so with probability bounded below there are at least \(k/8\) such hits. In any inspection order, conditional hit bounds and a union bound over choices give probability at most \(2^k(2d p_*)^{k/8}\) for this many hits. The column statement with the reverse sweep and nonbad \(y\) is identical. This proves exponentially small absolute row or column sums on these parts of \(L\), since factors \(\exp(O(k))\) from summing over \(J\) are affordable. Arbitrary absolute row and column sums are bounded by \(2^k\).

We also need a dispersion observation: from any \(x\), the absolute mass from each \(n^{-k}K_J\) to bad \textbf{distinct} \(y\) is \(\exp(-\Omega(k\log(n/k)))\). We include the details because the same unfavorable starts cannot just be treated as dispersed. For a fixed set \(E\) of end positions, occupation counts here satisfy the upper tail bound from independent cards with success probability \(|E|/n\), in the sense of the ordinary Chernoff estimate. For positive common test factors on end positions we have a product upper bound on expectation by single-card expectations. To see this for \(J\), work backwards through the layers: for tracked entries occupying both inputs of a switch, the product of the two test factors on outputs is at most the square of their arithmetic mean. Thus products propagate as upper bounds with the common factor averaged at each switch. At the start the single-path propagations are all equal, since one sweep sends each start to a uniform end. This gives the needed generating function domination, also after including the independent entries outside \(J\).

If \(y\) is bad, take the maximal witnessing cells (having the density and count conditions). They are disjoint; there are \(l\le k/2\) of them, with total volume at most \(k/p_*\), and at least \(\max(k/4,2l)\) of the entries in their union. There are at most \(\binom{2n}{l}\) choices. Using the just-proved count bound, the costs for a given \(l\) are bounded by
\[
 \binom{2n}{l}
 \left(\frac{C k}{n p_*}\right)^{\max(k/4,2l)}
\]
for an absolute \(C\). These sum to the claimed dispersion estimate; for example, in the combinatorial prefactor the logs are bounded by \(l\log(n/k)+O(k)\), absorbed with room in the other exponent.

Partition \(L\) according to bad starts and ends. On the bad-bad block the dispersion estimate is available, and on the other blocks we have the nonbad estimates using cancellation of isolated paths. The row/column sum bound for operator norm (using the coarse \(2^k\) for the opposite estimate if necessary) proves \eqref{ten:named:eq:2}.

We will also use the zero bound when \(\lambda\) has more than \(n/2\) rows. Indeed a layer is the projection averaging a product of pair groups. Only diagrams in the product of \(n/2\) two-box row diagrams can survive, so this follows from Pieri.

\subsection{An angle estimate on an incomplete rectangle}
We state the form of the estimate to be used in the induction. Consider \(n=b s\) slots in \(b\) rows, \(s\) columns. In the application both \(b,s\) are within a factor two of \(\sqrt n\). Let \(m\) slots be holes, arbitrary but fixed for the estimate, and use \(V\) as above on the \(u=n-m\) remaining slots, with \(q\le n\). Take:
\begin{enumerate}
\item projections \(E,F\) specifying \(G\)-type in each row, respectively each column, of non-hole slots (each row and column can have its own type),
\item a projection \(Q\) to one global \(G\)-type \((\mu,\nu)\), with \(S=\mathcal S(\mu,\nu)\).
\end{enumerate}

The row and column projections are for actions constant in the respective row or column. Write \(S_H,S_K\) for sums of the type entropies in the specification for the rows, columns, respectively. Then
\begin{equation}
 \| E Q F\|^2
 \le
 \min\left(1,\
   \exp\{S_H+S_K-S+C m+C(1+b+s)q^2\log(n+2)\}
          \right).
 \label{ten:named:eq:4}
\end{equation}
These type spaces and \(Q\) do not require coordinates to be arranged contiguously for rows or columns. Using graded slot order instead of ordinary order to regroup the factors also respects the type definitions, by block-diagonality for plus/minus. In particular \(E,F\) commute with \(Q\).

Here is a proof. For a row with \(v>0\) slots after omitting holes and type \(w=(\alpha,\gamma)\), take a block-diagonal density matrix \(A_i\) on \(V\) with block spectra given by \(\alpha/v,\gamma/v\). Then the projection to the type is bounded above in positive order by
\begin{equation}
 \exp\{\mathcal S(w)+C q^2\log(n+2)\}\ 
   \int_G (g A_i g^*)^{\otimes v}\,dg.
 \label{ten:named:eq:5}
\end{equation}
Here \(dg\) is normalized Haar measure. Indeed, on that type the integral is scalar, with the scalar given by taking trace on the \(G\)-irrep and dividing by its dimension. The density tensor action has the standard polynomial action and is the same on all copies; this holds also for singular densities by continuity. By highest weight this trace is at least \(\exp(-\mathcal S(w))\), proving \eqref{ten:named:eq:5}. Use arbitrary densities on fully holed rows. We have the analogous bounds with column densities \(A'_j\).

Let \(w_i,w'_j\) be the row and column types, including the trivial type on a fully holed line, and choose one sufficiently large absolute \(C\) in the local bounds. Put
\[
\begin{aligned}
 c_i&=\exp\{\mathcal S(w_i)+Cq^2\log(n+2)\},\\
 d_j&=\exp\{\mathcal S(w'_j)+Cq^2\log(n+2)\},\\
 c_H&=\prod_i c_i,\\
 c_K&=\prod_j d_j.
\end{aligned}
\]
For a fully holed line the projection and tensor power are scalar one, so the same coefficient, which is at least one, is valid. Let \(\mu_H,\mu_K\) be the products of normalized Haar measures on \(G\), with one factor per row and column, respectively. For orientation tuples \(\mathbf g=(g_i)_i\), \(\mathbf h=(h_j)_j\), put
\[
\begin{aligned}
 R_{\mathbf g}&=\bigotimes_{(i,j)\text{ occupied}}g_i A_i g_i^*,\\
 C_{\mathbf h}&=\bigotimes_{(i,j)\text{ occupied}}h_j A'_j h_j^*.
\end{aligned}
\]
The parity-preserving graded regrouping identifies these positive tensors with products of the local row or column tensor powers. Tensoring \eqref{ten:named:eq:5} within each family gives
\[
 E\preceq c_H\int R_{\mathbf g}\,d\mu_H,\qquad
 F\preceq c_K\int C_{\mathbf h}\,d\mu_K.
\]
Both measures are probabilities, and
\(c_H=\exp\{S_H+Cbq^2\log(n+2)\}\),
\(c_K=\exp\{S_K+Csq^2\log(n+2)\}\).
Apply Lemma~\ref{ten:positive-mixtures} with its positive operators equal to \(E,F\) and its middle operator equal to \(Q\). Since \(E,F\) are projections, it gives exactly
\[
 \|EQF\|\le\sqrt{c_Hc_K}\,
 \sup_{\mathbf g,\mathbf h}
 \|R_{\mathbf g}^{1/2}Q C_{\mathbf h}^{1/2}\|.
\]
The square roots of these tensor products are the tensor products of the individual density square roots. Thus, after the displayed factor \(\sqrt{c_Hc_K}\), it suffices to bound
\begin{equation}
 \left\| \left(\bigotimes_{\rm slots} A_i^{1/2}\right)
        Q
        \left(\bigotimes_{\rm slots} {A'}_j^{1/2}\right)\right\|
 \label{ten:named:eq:6}
\end{equation}
uniformly in orientations of the densities, absorbing the \(G\)-actions into the notation.

Let
\( M=b^{-1}\sum_i A_i\).
Use \(Y=(M+\rho I)^{1/2}\), \(\rho>0\), to flatten on the left in \eqref{ten:named:eq:6}. On the global type, \(Y^{\otimes u}\) has norm at most
\[
 (1+2q\rho)^{u/2}\exp(-S/2).
\]
This follows from the highest-weight density estimate in
Lemma~\ref{lem:signed-type-density}, after normalizing
$M+\rho I$ by its trace $1+2q\rho$. Let $\pi$ be the compact
carrier representation for $(\mu,\nu)$. The operator
$Y^{\otimes u}Q$ acts as $\pi(Y)$ on that carrier and as the
identity on its multiplicity space, and is zero on all other types.
Lemma~\ref{s:compact-coefficient-extraction} therefore gives
\[
 Y^{\otimes u}Q=\int_G f(g)g^{\otimes u}\,dg,
 \qquad
 \int_G|f(g)|\,dg
 \le(\dim\pi)(1+2q\rho)^{u/2}e^{-S/2}.
\]
The carrier dimension is at most
$\exp(Cq^2\log(n+2))$.

Insert \((Y^{-1})^{\otimes u}\) on the left of that integral to get \(Q\). We bound the remaining norm in the integrand, on the squared scale, by a product over cells of
\[
 \operatorname{tr}\big( Y^{-1} A_i Y^{-1}\, g A'_j g^*\big)
\]
by Hilbert-Schmidt on each cell. On the \textbf{full} rectangle these nonnegative quantities sum to at most \(n\), by flattening and the unit traces on the other side. Thus the product over the \(u\) cells is at most \((n/u)^u\le \exp(Cm)\), with empty products harmless. This proves the desired bound on \eqref{ten:named:eq:6} by letting \(\rho\) tend to zero, and gives \eqref{ten:named:eq:4}.

\subsection{The independent trace induction}
We now complete the separate proof using the preceding crowded-cell sparse estimate~\eqref{ten:named:eq:2}, the coefficient-integral proof of~\eqref{ten:named:eq:4}, and a three-regime closure with one fixed $C_0$.

For this proof, fix
\[
 a=.02,\qquad q_n=\lfloor n^a\rfloor,\qquad c=\delta/40.
\]
We take \eqref{ten:named:eq:2} available up to \(k=n^{1-c}\). We use constants and asymptotic estimates below independent of \(h\).

Proving it for \(\mu,\nu\) of lengths at most \(q_n\) at any sufficiently large \(n\) suffices at that size. Indeed, if either has more rows, trim each after \(q_n\) rows, sending \(w_*\) boxes in total to the named count. By containment and associativity, there is an occurrence with the trimmed \(\mu,\nu\) and some new partition of size \(m+w_*\). The entropy at the old size \(n-m\) exceeds the trimmed entropy by at least \(w_*\log q_n\), since trimmed row counts were each at most \((n-m)/q_n\), and separating out the trimmed counts only lowers the contribution from the kept counts upon renormalizing the latter. On the range in \eqref{ten:named:eq:8}, \(R_n\) can increase by at most \(w_* (\delta\log n+1)\). This proves the sufficiency for large \(n\).

We will work with sufficiently large dyadic \(n\), assuming \eqref{ten:named:eq:9} for smaller sizes. Consequently, all estimates using subproblem sizes are allowed to use \textbf{any} lengths for spin partitions there. Also, regardless of the right-hand side in \eqref{ten:named:eq:9}, in the range of \eqref{ten:named:eq:2} we can make \(W_n(\lambda)\le 1\) for nontrivial irreps, by taking \(h\) large depending on \(c\), since \(D_\lambda\le n^k\). And this inequality holds already above height \(n/2\). We explain the treatment of finite sizes after the induction estimates.

Use \(b\) rows of size \(s\), with \(b,s\) powers of two and within a factor two of \(\sqrt n\). Split a sweep at this scale. We have sweep operators supported on
\[
 H=(\mathfrak S_s)^b,\qquad K=(\mathfrak S_b)^s,
\]
using rows and columns, respectively; inside a group each acts by independent smaller sweeps. The order will not matter to the estimates. In the product representing \(B_n\), apply the Araki-Lieb-Thirring trace inequality for positive matrices (at power \(h\)). Thus a trace of singular values to \(2h\) on any representation is bounded above by
\begin{equation}
 \operatorname{Tr} A_H A_K,
 \label{ten:named:eq:10}
\end{equation}
where \(A_H,A_K\) are the corresponding tensor products within \(H,K\) of the positive group-algebra elements from the subproblem singular squares taken to \(h\) (using \(BB^*\) or \(B^*B\) as appropriate). In particular, on a factor of size \(s\), the regular traces by irrep of that positive element are \(W_s\). Positivity here and below can equivalently be defined by the regular representation.

Fix an occurrence in \eqref{ten:named:eq:9}, with lengths now at most \(q=q_n\), and write \(S=\mathcal S(\mu,\nu)\). Use a representation \(X\) with \(m\) slots taken by \textbf{distinct named} objects, one of each name, and the other slots carrying \(V\) as above with this \(q\). Names can be treated as additional plus objects for permuting the factors (no sign cost for them). Use the whole Hilbert space allowing the names in any slots and order, with each fixed placement carrying \(V^{\otimes(n-m)}\). Use \(Q\) selecting global type \((\mu,\nu)\); it fixes placements of names. It commutes with the position shuffle. This space with \(Q\) contains at least \(D_\tau\) copies of \([\lambda]\). This follows from inducing from the name and spin subsets and the description of spin types; in particular the names alone carry a regular representation on their slots. And
\begin{equation}
 \frac{D_\lambda}{D_\tau}
 \le \binom{n}{m}\exp(S)
 \label{ten:named:eq:11}
\end{equation}
by the dimension upper estimates for induction. It follows from \eqref{ten:named:eq:10} that \(W_n(\lambda)\) is bounded by \(D_\lambda/D_\tau\) times
\[
 \operatorname{Tr}_X Q A_H A_K .
\]

Here we analyze this last trace, which need not single out \(\tau\) or \(\lambda\). Expand in group elements. For a contribution with the \(m\) names in cells \(T\), \(|T|=m\), the product of a row and a column permutation must fix these individual cells, not just setwise. Hence each of the two permutations must fix these cells pointwise. The trace contribution for \(T\) is therefore \(m!\) times
\begin{equation}
 \operatorname{Tr} Q\  a_H a_K
 \label{ten:named:eq:12}
\end{equation}
on the remaining spins. In this expression:
\begin{itemize}
\item \(a_H\) is given by using in each row only the group-algebra coefficients on the subgroup fixing the named positions pointwise, acting by graded permutations on the other spins. Likewise \(a_K\) in the columns.
\item These give positive operators. Indeed coefficient restriction of a positive element to a subgroup is positive, for instance by taking a principal submatrix. Each commutes with global \(G\) and hence with \(Q\). Within its family of rows or columns, each can be treated in tensor product order by using graded regrouping of tensor factors. Such regroupings intertwine the position and spin-type descriptions above; the graded signs track only the two spin summands, which are preserved by the \(G\) actions. Thus, for example, its trace estimates on specified types in that family can use ordinary products of the estimates.
\end{itemize}

\subsection{Local stabilizer traces}
We give the local trace bound needed for \eqref{ten:named:eq:12}. In a row, write \(r=|T\cap\text{row}|\). For the part of its operator on a \(G\)-type \(w_i=(\alpha_i,\gamma_i)\) on the \(s-r\) slots, put \(S_i=\mathcal S(w_i)\). Then the trace of this positive part is bounded by
\begin{equation}
 \frac{1}{(s)_r}
 \exp\left\{ -(1-\beta_s) S_i + R_s(r)
             + C (q^2+\sqrt{s})\log(n+2) \right\},
 \label{ten:named:eq:13}
\end{equation}
where \((s)_r=s!/(s-r)!\).

For details, consider an \(\eta\) in \(\alpha_i*\gamma_i^{\rm t}\). Its block of the coefficient restriction only gets contributions from Fourier components in size \(s\) that extend \(\eta\) upon restricting to \(\mathfrak S_{s-r}\). This follows immediately also by viewing coefficients of Fourier components as matrix coefficients and restricting them. Keep just those components in size \(s\). A diagram extending \(\eta\) here occurs in \(\eta*\xi\) for some \(\xi\) on \(r\) slots, so the induction hypothesis bounds each regular trace from these components by \(\exp(\beta_s S_i+R_s(r))\). Coefficient restriction changes the total regular trace of the kept positive element by the group order ratio \(1/(s)_r\), using the identity coefficient. Consequently this bounds the trace on the \(\eta\)-block after multiplying by \(D_\eta\), and allowing also the number of kept partitions, at most \(\exp(C\sqrt{s}\log(n+2))\). But \(D_\eta\) itself has lower estimate \eqref{ten:named:eq:1}, and the multiplicities to be included on the spin representation, as well as enumerating its \(\eta\) here, are absorbed in the displayed error. This proves \eqref{ten:named:eq:13}. This argument pays only for diagrams \emph{which can go through the type and the \(r\) other slots}. In particular, the small named-object budget in \eqref{ten:named:eq:9} applies here even if \(r\) is just a local count and the global count is much larger.

Within \eqref{ten:named:eq:12}, take a type specification \(E\) for the rows and \(F\) for the columns as in \eqref{ten:named:eq:4}. The operators \(a_H,a_K\) break into positive parts according to their own families' specifications, since they commute with the family actions of \(G\). Also, \(Q\) commutes with the types of either family. By \eqref{ten:named:eq:13}, the traces of the two corresponding parts are bounded on the exponential scale using
\[
 -(1-\bar\beta)(S_H+S_K)
 \quad\text{in their product},\qquad
 \bar\beta=\max(\beta_s,\beta_b),
\]
besides the falling factorials, \(R\) terms and errors, where \(S_H,S_K\) are sums of the \(S_i\)'s and the analogous column entropies. One more factor in bounding \eqref{ten:named:eq:12} is at most
\(\|E Q F\|^2\): one can use eigenvectors of the positive parts and insert \(Q\) again since the operators commute with it. Combining \eqref{ten:named:eq:4} with the traces, we can replace the exponent just shown by
\begin{equation}
 -(1-\bar\beta) S + C m + C(1+b+s)q^2\log(n+2).
 \label{ten:named:eq:14}
\end{equation}
Indeed if the sum entropies reach \(S\) the trace term suffices, and otherwise the squared angle makes up the difference with rate 1.

\subsection{Counting the named positions}
All errors here other than \(O(m)\), and enumerations of the type specifications, cost \(O(n^{.85})\) in the exponent for large \(n\). For example \eqref{ten:named:eq:13} is used on \(O(\sqrt n)\) fibers, and even using the partition count per fiber on this scale is affordable. Thus with \(r_i=|T\cap\text{row }i|\) and \(t_j=|T\cap\text{column }j|\), equations \eqref{ten:named:eq:11}--\eqref{ten:named:eq:14} give
\begin{equation}
\begin{split}
 W_n(\lambda)\ \le\
 & \binom{n}{m}
 \exp\{\,\bar\beta S+C_1 m+C_1 n^{.85}\,\}\\
 &{}\times m!\sum_{T:\, |T|=m}
 \frac{\exp\{ \sum_i R_s(r_i)+\sum_j R_b(t_j)\}}
      {\prod_i(s)_{r_i}\prod_j(b)_{t_j}},
\end{split}
 \label{ten:named:eq:15}
\end{equation}
with an absolute \(C_1\). We emphasize that this constant does not use \(C_0\) in \eqref{ten:named:eq:8}. The elementary falling factorial bounds
\((s)_r\ge s^r\exp(-C r)\) and
\(m!\binom{n}{m}\le n^m\)
can be used to treat the last line simply as an expectation over uniform \(T\), with at most \(\exp(Cm)\) extra factor. The \(\binom{n}{m}\) on the first line, which could still be expensive, is why we propagate with the particular \(R\) in \eqref{ten:named:eq:8}.

\subsection{Closing the budget in three regimes}
We verify \eqref{ten:named:eq:9} from \eqref{ten:named:eq:15}. First suppose
\[
 m\ge n^{1-\delta/4}.
\]
If \(n\) is sufficiently large (permitting dependence on \(C_0\)), then \(\delta\log s-\log(n/m)-C_0>0\), and analogously for \(b\). Hence for instance
\[
 R_s(r_i)\le
 r_i\left(\delta\log s-\log(n/m)-C_0+
                 [\log(r_i/(m/b))]_+\right).
\]
The extra positive-part terms for \textbf{both} families together can be paid for in \eqref{ten:named:eq:15} with cost \(\exp(Cm)\). To see this, compare uniform \(T\) with iid cells before requiring distinctness, at likelihood cost bounded by \(\exp(Cm)\). Row and column counts of these iid cells arise independently. The row profile probability is bounded by
\(\exp(-\sum r_i\log(r_i/(m/b)))\), and likewise for columns. The negative-log entries discarded in using the positive part cost \(O(m)\) in each family (by \(z\log(1/z)\le C\) applied to ratios). And enumerating profiles costs \(\exp(Cm)\) here. Thus the budget terms on the last line cost in the exponent at most
\[
 m(\delta\log n-2\log(n/m)-2C_0)+C m,
\]
with constants independent of \(C_0\). Using \(\log\binom{n}{m}\le m\log(n/m)+m\), \eqref{ten:named:eq:15} now fits into \eqref{ten:named:eq:9}, including its \(R_n(m)\), by fixing \(C_0\) sufficiently large and taking \(n\) sufficiently large. The size errors \(C_1 n^{.85}\) in this case can be treated just as \(O(m)\).

Next, if \(m<n^{1-\delta/4}\) but \(S\ge n^{1-\delta/6}\), we can use crude cost \(O(m\log n)\) for the budget and binomial terms together. All overheads then fit in \((\beta_n-\bar\beta)S\); the coefficient difference is, in particular, at least a constant times \(1/(1+\log n)\) for large \(n\).

In the remaining case (\(m<n^{1-\delta/4}\) and \(S<n^{1-\delta/6}\)), one of the entries of \((\mu,\nu)\) differs from \(n-m\) by \(O(n^{1-\delta/6})\), just by entropy. Thus because \(\lambda\) contains both \(\mu\) and \(\nu^{\rm t}\), either it is above the height cutoff, or its level \(n-\lambda_1\) is in the sparse range used for \eqref{ten:named:eq:2}, for large \(n\). This gives \eqref{ten:named:eq:9} by the sparse bound in \eqref{ten:named:eq:7} or the zero bound, the trivial irrep also satisfying \eqref{ten:named:eq:9}.

This closes the trace induction. To make the choices with fixed \(h\) explicit as regards dependencies, choose \(C_0\) to meet the absolute constant requirements in the verification, use the arguments just given (and the trimming observation) above some absolute finite size, and take \(h\) to meet the requirement for the sparse bound and to give \eqref{ten:named:eq:9} on all needed finite sizes. The comparisons above for sufficiently large \(n\) can choose their thresholds independently of further increases in \(h\). For the finite sizes a fixed such \(h\) exists: in any nontrivial irrep the sweep has strict contraction. Indeed, it is a product of orthogonal projections, and equality in norm would require a common fixed vector. The pair interchanges in the coordinate lines generate the symmetric group. We include subproblems of one slot as trivial. These choices prove \eqref{ten:named:eq:9} with absolute constants.

\subsection{Dimension savings and full-deck mixing}
We show why the propagated budget suffices to mix, in particular that it handles shapes with many rows and also many columns. Continue to use \(c=\delta/40\). Small levels and diagrams above the height cutoff are already handled. For any remaining diagram \(\lambda\), take the following occurrence for use in \eqref{ten:named:eq:9}:
\begin{itemize}
\item take the first up to \(q_n\) rows as \(\mu\);
\item below them take the diagram within the first \(q_n\) columns, transposed, as \(\nu\);
\item the remaining southeast boxes, number \(m\), can be assigned a partition factor on \(m\) slots.
\end{itemize}

More precisely, let \(\eta\) be the hook portion kept here. We have \(\eta\) in \(\mu*\nu^{\rm t}\); this uses adjoining the diagram \(\nu^{\rm t}\) under upper rows whose lengths cover its width (or follows directly by the Littlewood-Richardson rule). And since \(\lambda\) contains \(\eta\), some \(m\)-slot extension suffices. In particular branching and \eqref{ten:named:eq:1} give
\[
 \log D_\lambda\ge S-C q_n^2\log(n+2),
 \qquad S=\mathcal S(\mu,\nu).
\]
If \(m\ge n^{1-2\delta}\), we also have
\(\log D_\lambda\ge (a/2)m\log n\)
for large \(n\). Indeed the remaining boxes alone form, upon translation, a diagram of size \(m\) with every row and column length at most \(n/q_n\). We have at least its dimension in the tableau count for \(\lambda\), just by filling it after the hook part. The estimate follows by the hook and factorial formulas.

Here at least one of \(m,S\) is \(\ge n^{1-2\delta}\) for large \(n\). Otherwise we would again be above the height cutoff or in the sparse range (recall that the remaining irreps under discussion have level greater than \(n^{1-c}\)). Thus \(D_\lambda\) is at least \(\exp(n^{.9})\) for large \(n\), and all additive errors in these dimension estimates are harmless. Since \(R_n(m)=0\) unless \(m\) is large enough for the second estimate, and \(\delta\) is tiny compared with \(a\), we see that
\[
 W_n(\lambda)\le D_\lambda^{1/2}
\]
on all these remaining diagrams, for sufficiently large \(n\). In particular the sweep on them has norm at most \(D_\lambda^{-1/(4h)}\).

Use now an absolute constant number \(L\) of sweeps, repeating in the same order. In squared \(L^2\) distance of density from uniform on the permutation group, Plancherel uses, for each nontrivial irrep, the squared Hilbert-Schmidt norm times \(D_\lambda\), and thus we can bound by summing
\[
 D_\lambda^2\,\|B_n([\lambda])\|^{2L}.
\]
On the large dimensions just treated, this gives \(o(1)\) in total for \(L\) sufficiently large, even allowing the partition count \(\exp(O(\sqrt n\log(n+2)))\). On the sparse range \(1\le k\le n^{1-c}\), use \eqref{ten:named:eq:2}, \(D_\lambda\le n^k\), and for example \(2^{O(k)}\) partitions per level. This sum also tends to zero for large absolute \(L\). This proves in particular the required TV upper bound by \(O(d)\) steps, from every fixed initial order.
